\documentclass[../main]{subfiles}
\begin{document}

\chapter{Ensayos Eléctricos en Motores Asincrónicos}

\section{Introducción a los Motores Asincrónicos}

Los motores asincrónicos, también conocidos como motores de inducción, son ampliamente utilizados en aplicaciones industriales debido a su robustez y simplicidad constructiva. Funcionan basados en el principio de inducción electromagnética, donde el campo magnético giratorio generado en el estator induce una corriente en el rotor, produciendo así el par motor.

\section{Ensayos Eléctricos Característicos}

Para comprender el funcionamiento y la eficiencia de un motor asincrónico, se realizan varios ensayos eléctricos. Estos ensayos permiten analizar distintos fenómenos energéticos que ocurren durante el funcionamiento del motor. A continuación, se describen algunos ensayos importantes.

\subsection{Ensayo de Vacío}

El ensayo de vacío se realiza con el motor desconectado de cualquier carga mecánica y a su tensión nominal. Este ensayo permite determinar las pérdidas en el núcleo y las pérdidas por fricción y ventilación. La potencia absorbida en este ensayo se utiliza para calcular el rendimiento del motor.

\info{Fórmula de Potencia de Entrada en Ensayo de Vacío}{
La potencia de entrada \( P_0 \) en el ensayo de vacío se calcula como:
\begin{equation}
P_0 = V \cdot I_0 \cdot \cos(\phi_0)
\end{equation}
Donde:
\begin{itemize}
    \item $V$: Voltaje aplicado [V]
    \item $I_0$: Corriente de vacío [A]
    \item $\cos(\phi_0)$: Factor de potencia en vacío
\end{itemize}
}

\subsection{Ensayo de Carga}

El ensayo de carga se realiza conectando el motor a una carga mecánica conocida. Este ensayo permite evaluar el comportamiento del motor bajo condiciones de operación normales y determinar el par motor y la eficiencia a diferentes niveles de carga.

\info{Fórmula del Par Motor}{
El par motor \( T \) se calcula utilizando la siguiente fórmula:
\begin{equation}
T = \frac{P_{out}}{\omega}
\end{equation}
Donde:
\begin{itemize}
    \item $P_{out}$: Potencia de salida [W]
    \item $\omega$: Velocidad angular [rad/s]
\end{itemize}
}

\subsection{Ensayo de Rotor Bloqueado}

En el ensayo de rotor bloqueado, el rotor se inmoviliza y se aplica un voltaje reducido. Este ensayo permite determinar las pérdidas en el rotor y el circuito equivalente del motor. Es útil para calcular la resistencia y reactancia del rotor.

\info{Fórmula de la Resistencia del Rotor}{
La resistencia del rotor \( R_r \) se determina mediante:
\begin{equation}
R_r = \frac{V_{bl}}{I_{bl}}
\end{equation}
Donde:
\begin{itemize}
    \item $V_{bl}$: Voltaje aplicado en rotor bloqueado [V]
    \item $I_{bl}$: Corriente en rotor bloqueado [A]
\end{itemize}
}

\section{Ejemplo de Cálculo}

\ejemplo{Ejemplo de Ensayo de Vacío}{
Calcular la potencia de entrada en un ensayo de vacío donde se aplica un voltaje de 230V, la corriente de vacío es 2A y el factor de potencia en vacío es 0.1.

Datos:
\begin{itemize}
    \item $V = 230 \, V$
    \item $I_0 = 2 \, A$
    \item $\cos(\phi_0) = 0.1$
\end{itemize}

\begin{equation}
P_0 = 230 \, V \cdot 2 \, A \cdot 0.1 = 46 \, W
\end{equation}
}

\actividad{Responder el siguiente cuestionario}
\begin{enumerate}
    \item ¿Qué es un motor asincrónico y en qué principio se basa su funcionamiento?
    \item ¿Qué tipo de pérdidas se determinan en el ensayo de vacío?
    \item ¿Cómo se calcula la potencia de entrada en el ensayo de vacío?
    \item ¿Cuál es el propósito del ensayo de carga en un motor asincrónico?
    \item ¿Qué información proporciona el ensayo de rotor bloqueado?
    \item ¿Qué significa el factor de potencia en el contexto de motores eléctricos?
    \item ¿Por qué es importante conocer la resistencia del rotor?
    \item ¿Qué parámetros se miden en el ensayo de vacío?
    \item ¿Cómo se relaciona la velocidad angular con el par motor?
    \item ¿Qué factores influyen en la eficiencia de un motor asincrónico?
\end{enumerate}


\actividad{Resolver los siguientes problemas}
\begin{enumerate}
    \item Calcular la potencia de entrada en un ensayo de vacío para un motor al que se le aplica un voltaje de 220V, la corriente de vacío es de 1.5A y el factor de potencia es 0.12.
    \item Determinar el par motor generado si la potencia de salida es de 500W y la velocidad angular es de 157 rad/s.
    \item Hallar la resistencia del rotor si en un ensayo de rotor bloqueado se aplica un voltaje de 50V y la corriente medida es de 10A.
    \item Calcular la corriente de vacío si se sabe que la potencia de entrada es de 100W, el voltaje aplicado es 230V y el factor de potencia es 0.1.
    \item Determinar la velocidad angular de un motor si se sabe que el par motor es de 20Nm y la potencia de salida es de 1000W.
    \item Calcular la corriente en rotor bloqueado si la resistencia del rotor es 5Ω y el voltaje aplicado es de 25V.
    \item Hallar el voltaje aplicado en un ensayo de vacío si la potencia de entrada es de 92W, la corriente de vacío es de 1A y el factor de potencia es 0.4.
    \item Determinar el factor de potencia en vacío si la potencia de entrada es de 36W, el voltaje aplicado es 120V y la corriente de vacío es de 1.5A.
    \item Calcular la potencia de salida si el par motor es de 30Nm y la velocidad angular es de 314 rad/s.
    \item Hallar la resistencia del rotor si se sabe que la corriente en rotor bloqueado es de 8A y el voltaje aplicado es de 40V.
\end{enumerate}

\end{document}

